% from ega4-1-fr.tex
\oldpage[IV-1]{234}

\begin{proposition}[1.5.3]
\label{IV.1.5.3-fr}
Soit $f:X\to Y$ un morphisme de type fini~; si $Y$ est noethérien, il en est
de même de $X$.
\end{proposition}

Pour la démonstration, voir (I, 6.3.7).

\begin{proposition}[1.5.4]
\label{IV.1.5.4-fr}
\begin{enumerate}
  \item[(i)] Soit $j:X\to Y$ un morphisme d'immersion. Si $j$ est
  quasi-compact (ce qui est le cas si $j$ est une immersion fermée, ou si
  l'espace sous-jacent à $X$ est noethérien, ou si l'espace sous-jacent à $Y$
  est localement noethérien), $j$ est de type fini.
  \item[(ii)] Le composé de deux morphismes de type fini est de type fini.
  \item[(iii)] Si $f:X\to Y$ est un $S$-morphisme de type fini,
  $f_{(S')}:X_{(S')}\to Y_{(S')}$ est de type fini pour toute extension
  $S'\to S$ du préschéma de base.
  \item[(iv)] Si $f:X\to X'$ et $g:Y\to Y'$ sont deux $S$-morphismes de type
  fini, $f\times_Sg$ est de type fini.
  \item[(v)] Si le composé $g\circ f$ de deux morphismes est de type fini et
  si $g$ est quasi-séparé ou si $X$ est noethérien, $f$ est de type fini.
  \item[(vi)] Si un morphisme $f$ est de type fini, il en est de même de
  $f_{\mathrm{red}}$.
\end{enumerate}
\end{proposition}

Cela résulte aussitôt (sauf le cas (v) où $X$ est noethérien, démontré dans
(I, 6.3.6)) des définitions et de (1.1.2), (1.2.4) et (1.3.4).

\begin{corollary}[1.5.5]
\label{IV.1.5.5-fr}
Soit $f:X\to Y$ un morphisme de type fini. Pour tout morphisme $Y'\to Y$ tel
que $Y'$ soit noethérien, $X\times_YY'$ est noethérien.
\end{corollary}

Cela résulte de (1.5.4, (iii)) et de (1.5.3).

\begin{corollary}[1.5.6]
\label{IV.1.5.6-fr}
Soit $X$ un préschéma de type fini sur un préschéma localement noethérien $S$.
Alors tout $S$-morphisme $f:X\to Y$ est de type fini.
\end{corollary}

Pour la démonstration, voir (I, 6.3.9).

\begin{proposition}[1.5.7]
\label{IV.1.5.7-fr}
Soit $\varphi:A\to B$ un homomorphisme d'anneaux. Pour que le morphisme
correspondant $f:\operatorname{Spec}(B)\to\operatorname{Spec}(A)$ soit de type fini, il faut et il suffit
que $\varphi$ fasse de $B$ une $A$-algèbre de type fini.
\end{proposition}

Pour la démonstration, voir (I, 6.3.3).

\subsection{Morphismes de présentation finie}
\label{subsection:IV.1.6-fr}

\begin{definition}[1.6.1]
\label{IV.1.6.1-fr}
Soient $X$, $Y$ deux préschémas. On dit qu'un morphisme $f:X\to Y$ est
\emph{de présentation finie} s'il vérifie les trois conditions suivantes~:
\begin{enumerate}
  \item[(i)] $f$ est localement de présentation finie~;
  \item[(ii)] $f$ est quasi-compact (ce qui, lorsque (i) est vérifiée,
  équivaut à dire que $f$ est de type fini (1.5.2))~;
  \item[(iii)] $f$ est quasi-séparé.
\end{enumerate}
On dit alors aussi que $X$ est \emph{de présentation finie sur $Y$}, ou un
\emph{$Y$-préschéma de présentation finie}.
\end{definition}

La condition (iii) est automatiquement vérifiée si $f$ est séparé, ou si $X$
est localement noethérien (1.2.8)~; lorsque $Y$ est localement noethérien, il
revient au même de dire que $f$ est \emph{de présentation finie} ou que $f$
est \emph{de type fini} (cette dernière condition entraînant que $X$ est
localement noethérien (1.3.3)).

